\lysh{20653_SOLUTION}{1}
α) Έστω $t \sec$ ο χρόνος που μεσολάβησε από τη στιγμή της εκπυρσοκρότησης μέχρι τη στιγμή που την άκουσε ο παρατηρητής $M_2$. Τότε ο αντίστοιχος χρόνος για τον $M_1$ θα είναι $t+4 \sec$. Άρα $(PM_2) = 340 \cdot t$ και $(PM_1) = 340 \cdot (t+4) = 340 \cdot t + 1360$.
Ώστε $(PM_1) - (PM_2) = 1360 \,m$.

β) Γνωρίζουμε ότι όλα τα σημεία $P$ που ικανοποιούν τη σχέση $|(PM_1) - (PM_2)| =$ σταθερή, ανήκουν σε δύο κλάδους υπερβολής με εστίες τα σταθερά σημεία $M_1$ και $M_2$. Άρα η θέση $P$ του αγνοούμενου θα ανήκει σε έναν κλάδο υπερβολής με εστίες τα $M_1$ και $M_2$ και προφανώς σε αυτόν που βρίσκεται πιο κοντά στον παρατηρητή $M_2$.

γ) Η ζητούμενη εξίσωση είναι της μορφής $\frac{x^2}{\alpha^2} - \frac{y^2}{\beta^2} = 1$. Εδώ είναι $2\alpha = 1360$, οπότε $\alpha = 680$. Η απόσταση $(M_1M_2)$ είναι $2\gamma$, άρα $\gamma = 1378 : 2 = 689$.
Όμως
\[
\beta^2 = \gamma^2 - \alpha^2 = 689^2 - 680^2 = (689 + 680)(689 - 680) = 1369 \cdot 9 = 37^2 \cdot 3^2 = 111^2.
\]

% TODO: TikZ σχήμα
\end{lysh}